% Regression test: the editing rules that drop edges. \proofgraphignore takes
% the labels of \proofgraphedge, in that order, the result first, and drops
% exactly that one dependency; \proofgraphignorecite does the same for a
% citation node, per (key, note) pair, so ignoring the noted citation leaves the
% plain one standing and ignoring the last dependency on a work removes its node
% too.
% A rule written the other way round matches nothing and must leave the graph
% untouched (it is reported as a warning, which the .dot does not show).
% Uses only amsthm; single-pass and bibtex-free (citelabel=key, manual \bibitem)
% for determinism.
\documentclass{article}
\usepackage{amsthm}
\usepackage[cite,citelabel=key]{proofgraph}

\newtheorem{theorem}{Theorem}
\newtheorem{lemma}[theorem]{Lemma}

\proofgraphignore{thm:m}{lem:a}
\proofgraphignorecite[Thm~3]{thm:m}{bk} % only the noted node of bk

\begin{document}

\begin{lemma}\label{lem:a}A.\end{lemma}
\begin{lemma}\label{lem:b}B.\end{lemma}
\begin{theorem}\label{thm:m}Main.\end{theorem}
\begin{proof}
  By \ref{lem:a} and \ref{lem:b}, \cite{bk} and \cite[Thm~3]{bk},
  and \cite{art}.
\end{proof}

\begin{theorem}\label{thm:n}Other.\end{theorem}
\proofgraphedge{thm:n}{lem:b}

\proofgraphignorecite{thm:m}{art}   % last edge to art: its node goes as well
\proofgraphignore{lem:b}{thm:n}     % the wrong way round: no effect

\begin{thebibliography}{9}
\bibitem{bk}A book.
\bibitem{art}An article.
\end{thebibliography}
\end{document}
